% Problem record op_09b9fa91a1ac1a76. This TeX file is derived from
% database/problems_json/op_09b9fa91a1ac1a76.json by scripts/sync-tex.mjs; edit the JSON record
% and rerun the script rather than editing this file.
\section{Asymptotic metrology with quantum-controlled causal order}
\label{sec:op_09b9fa91a1ac1a76}

\paragraph{Problem.}
Does quantum control of causal order yield a persistent asymptotic
metrological advantage over parallel access for some regular,
finite-dimensional channel family? Let
$\Lambda_\theta:\mathcal L(A)\to\mathcal L(B)$ be a smooth one-parameter
family on an open interval. A parameter value $\theta$ is regular here if
the Choi rank is constant in a neighborhood of $\theta$; the family then
admits a differentiable minimal Kraus representation locally.
Let $\mathcal F_{\mathrm{PAR}}^{(N)}(\theta)$ and
$\mathcal F_{\mathrm{QCQC}}^{(N)}(\theta)$ be the largest output-state
symmetric-logarithmic-derivative quantum Fisher information attainable from
$N$ black-box uses by, respectively, a parallel strategy and a quantum
circuit with quantum control of causal order (QC-QC). Both classes allow
arbitrary finite noiseless ancillary systems. In a QC-QC strategy, a
coherent control can dynamically select which unused black-box call occurs
next; this is more general than coherently superposing fixed causal
orders. At every regular parameter value with
$\mathcal F_{\mathrm{PAR}}^{(1)}(\theta)>0$, decide whether every such family
satisfies
\begin{equation}
  \limsup_{N\to\infty}
  \frac{\mathcal F_{\mathrm{QCQC}}^{(N)}(\theta)}
       {\mathcal F_{\mathrm{PAR}}^{(N)}(\theta)}
  =1.
  \label{eq:p61-asymptotic-qcqc-ratio}
\end{equation}
Equivalently, either prove Eq.~\eqref{eq:p61-asymptotic-qcqc-ratio} or
construct a noisy family in this regular domain and physical QC-QC
strategies with a persistent constant-factor advantage or a larger
scaling exponent.

\subsection*{Status}

Solved

\subsection*{Source}

Mothe, Branciard, and Abbott explicitly ask whether the finite-use QC-QC
advantage that they establish can persist asymptotically
\sourcecite{ref:p61-mothe}{MBA24}.

\subsection*{Progress}

\begin{itemize}
  \item Kurdzia\l ek, G\'orecki, Albarelli, and Demkowicz-Dobrza\'nski prove
  that adaptive strategies and coherent superpositions of fixed causal
  orders are asymptotically equivalent to parallel strategies for repeated
  estimation of a finite-dimensional channel
  \sourcecite{ref:p61-kurdzialek}{KGAD23}.  Their recursion does not cover
  the dynamically controlled QC-QC class in
  Eq.~\eqref{eq:p61-asymptotic-qcqc-ratio}.

  \item Mothe, Branciard, and Abbott exhibit noisy channel families for which
  QC-QC strictly outperforms parallel, sequential, and causal-superposition
  strategies at three uses \sourcecite{ref:p61-mothe}{MBA24}.  This proves a
  finite-resource separation but neither a nonunit asymptotic ratio nor a
  different scaling exponent.

  \item Abbott, Mhalla, and Pocreau show that QC-QC gives no query advantage
  for Boolean-function tasks formed by repeated access to one unitary
  \sourcecite{ref:p61-abbott}{AMP24}.  The restriction to unitary queries
  does not address the noisy metrological families relevant to
  Eq.~\eqref{eq:p61-asymptotic-qcqc-ratio}.

  \item Salzger and Vilasini show that higher-order processes satisfying
  their spacetime, Acting Once, and Local Order assumptions reduce
  operationally to QC-QC processes
  \sourcecite{ref:p61-salzger}{SV26}.  This identifies QC-QC as the physical
  target class under those assumptions but supplies no asymptotic Fisher
  information bound.

  \item Wei, Li, and Pang's preprint establishes an upper bound for general
  indefinite-causal-order processes in Theorem~3 and Supplementary
  Section~V \sourcecite{ref:p61-wei}{WLP26}. For a differentiable minimal
  Kraus representation, write $\alpha=\sum_i\dot K_i^\dagger\dot K_i$ and
  $\beta=\sum_i\dot K_i^\dagger K_i$. Combining that upper bound with the
  parallel attainability results of Zhou and Jiang, Theorems~2 and~3
  \sourcecite{ref:p61-zhou}{ZJ21}, gives
  \begin{equation}
    \lim_{N\to\infty}\frac{\mathcal F_{\mathrm{Gen}}^{(N)}(\theta)}{N^s}
    =\lim_{N\to\infty}\frac{\mathcal F_{\mathrm{PAR}}^{(N)}(\theta)}{N^s}
    =c>0,
    \label{eq:p61-leading-coefficient}
  \end{equation}
  where $s=1$ and $c=4\min_{\beta=0}\lVert\alpha\rVert$ if a Kraus gauge
  with $\beta=0$ exists; otherwise $s=2$ and
  $c=4\min\lVert\beta\rVert^2$. All minima are over Kraus gauges at the
  parameter value in question. In the linear regime, product inputs give
  $c\ge\mathcal F_{\mathrm{PAR}}^{(1)}(\theta)>0$; in the quadratic regime,
  positivity follows from the nonzero minimum defining that regime.
  Since $\mathrm{PAR}\subseteq\mathrm{QCQC}\subseteq\mathrm{Gen}$,
  Eq.~\eqref{eq:p61-leading-coefficient} proves
  Eq.~\eqref{eq:p61-asymptotic-qcqc-ratio} by squeezing the ratio.
\end{itemize}

\subsection*{References}

\begin{enumerate}[label={},leftmargin=4.8em,labelsep=0.5em]
  \footnotesize
  \item[\textup{[MBA24]}]\label{ref:p61-mothe}
  R. Mothe, C. Branciard, and A. A. Abbott,
  ``Reassessing the Advantage of Indefinite Causal Orders for Quantum
  Metrology,'' \emph{Physical Review A} \textbf{109}, 062435 (2024).
  \href{https://doi.org/10.1103/PhysRevA.109.062435}{doi:10.1103/PhysRevA.109.062435};
  \href{https://arxiv.org/abs/2312.12172}{arXiv:2312.12172}.

  \item[\textup{[KGAD23]}]\label{ref:p61-kurdzialek}
  S. Kurdzia\l ek, W. G\'orecki, F. Albarelli, and
  R. Demkowicz-Dobrza\'nski,
  ``Using Adaptiveness and Causal Superpositions Against Noise in Quantum
  Metrology,'' \emph{Physical Review Letters} \textbf{131}, 090801 (2023).
  \href{https://doi.org/10.1103/PhysRevLett.131.090801}{doi:10.1103/PhysRevLett.131.090801};
  \href{https://arxiv.org/abs/2212.08106}{arXiv:2212.08106}.

  \item[\textup{[AMP24]}]\label{ref:p61-abbott}
  A. A. Abbott, M. Mhalla, and P. Pocreau,
  ``Quantum Query Complexity of Boolean Functions under Indefinite Causal
  Order,'' \emph{Physical Review Research} \textbf{6}, L032020 (2024).
  \href{https://doi.org/10.1103/PhysRevResearch.6.L032020}{doi:10.1103/PhysRevResearch.6.L032020};
  \href{https://arxiv.org/abs/2307.10285}{arXiv:2307.10285}.

  \item[\textup{[SV26]}]\label{ref:p61-salzger}
  M. Salzger and V. Vilasini,
  ``Higher-Order Quantum Processes Respecting Closed Labs in a Spacetime
  Have Quantum-Controlled Causal Order,'' arXiv:2605.08351 (2026).
  \href{https://arxiv.org/abs/2605.08351}{arXiv:2605.08351}.

  \item[\textup{[WLP26]}]\label{ref:p61-wei}
  W. Wei, Y. Li, and S. Pang,
  ``Fundamental Limits of Quantum Metrology Beyond Fixed Causal Order,''
  arXiv:2609.05355v1 (2026).
  \href{https://arxiv.org/abs/2609.05355v1}{arXiv:2609.05355v1}.

  \item[\textup{[ZJ21]}]\label{ref:p61-zhou}
  S. Zhou and L. Jiang,
  ``Asymptotic Theory of Quantum Channel Estimation,''
  \emph{PRX Quantum} \textbf{2}, 010343 (2021).
  \href{https://doi.org/10.1103/PRXQuantum.2.010343}{doi:10.1103/PRXQuantum.2.010343};
  \href{https://arxiv.org/abs/2003.10559}{arXiv:2003.10559}.
\end{enumerate}

\subsection*{Comment}

The asymptotic-ratio question in the explicitly regular,
ancilla-assisted domain above is resolved by Wei, Li, and Pang's
Theorem~3, together with the parallel attainability theorem. The resolving
Wei--Li--Pang result is a preprint, arXiv:2609.05355v1, rather than a
peer-reviewed publication. Rank-changing parameter values are outside the
regular domain defined in the statement. Finite-use advantages and
subleading corrections can persist without changing the limiting ratio
in Eq.~\eqref{eq:p61-asymptotic-qcqc-ratio}.

\subsection*{Field}

Quantum metrology; Quantum Resource Theory

\subsection*{Topic}

Quantum estimation; Indefinite causal order

\subsection*{ID}

\texttt{op\_09b9fa91a1ac1a76}
